\documentclass[10pt,a4paper]{article}
\pdfinfoomitdate=1
\pdftrailerid{}
\pdfsuppressptexinfo=15
\usepackage[T1]{fontenc}
\usepackage[utf8]{inputenc}
\usepackage{lmodern}
\usepackage{microtype}
\usepackage[a4paper,margin=22mm]{geometry}
\usepackage{booktabs,tabularx,array}
\usepackage{enumitem}
\usepackage{xcolor}
\usepackage{hyperref}
\usepackage{xurl}
\usepackage{fancyhdr}
\definecolor{navy}{HTML}{17324D}
\definecolor{teal}{HTML}{1B728C}
\definecolor{soft}{HTML}{F3F6F8}
\definecolor{warn}{HTML}{8A4B08}
\definecolor{greenish}{HTML}{2E6B57}
\hypersetup{colorlinks=true,linkcolor=navy,urlcolor=teal,
  pdftitle={Siamese Embedding Compression Research Program v0.3.2 Status Addendum},
  pdfauthor={Guillaume Harbonnier},
  pdfsubject={Phase A construct-validity acceptance with limitations and prospective closure of B1-B3}}
\setlength{\parindent}{0pt}
\setlength{\parskip}{5pt}
\setlist[itemize,enumerate]{nosep,leftmargin=*}
\pagestyle{fancy}
\fancyhf{}
\fancyhead[L]{Siamese Embedding Compression Research Program}
\fancyhead[R]{v0.3.2 status addendum}
\fancyfoot[C]{\thepage}
\renewcommand{\headrulewidth}{0.4pt}
\newcommand{\code}[1]{\texttt{#1}}
\title{\textbf{Siamese Embedding Compression Research Program}\\[4pt]
\large Status Addendum: construct validity accepted with prospective blockers}
\author{Guillaume Harbonnier\\
\small Independent exploratory research\\
\small \url{https://github.com/gharbonnier78/siamese-embedding-compression-lab}}
\date{Version 0.3.2 --- 16 September 2026}
\begin{document}
\maketitle
\thispagestyle{empty}
\begin{center}
\fcolorbox{warn}{soft}{\parbox{0.92\linewidth}{
\textbf{Scope.} This addendum updates v0.3.1 after the independent Phase A construct-validity review. Construct validity
is accepted with limitations on the exact reviewed head \code{b5720dd7...}; three prospective blockers B1--B3 must be
closed before platform materialization. This addendum opens no protected Study 1B outcome, launches no S4N3, selects no
target platform and does not authorize canonical Phase A execution.
}}
\end{center}

\section{Independent construct-validity verdict}

The independent review of the N7/N8-hardened package bound to exact head \code{b5720dd7...} returned:

\begin{quote}
\code{VERDICT: ACCEPT\_WITH\_LIMITATIONS}\\
\code{CONSTRUCT\_VALIDITY\_ACCEPTED: yes}\\
\code{REQUIRES\_PROTOCOL\_CHANGE\_BEFORE\_PLATFORM\_MATERIALIZATION: yes}\\
\code{CANONICAL\_EXECUTION\_CURRENTLY\_ADMISSIBLE: no}
\end{quote}

All ten construct-validity dimensions were accepted: intervention isolation, fixed-work dense-search semantics,
synthetic-payload adequacy for the bounded engineering characterization, G0/G2 scenario scope, endpoint alignment,
measurement-plane validity, confound control, generalization limits, characterization-versus-qualification and the
interpretation firewall.

Transport hardening was also independently checked: 13/13 repository Git blobs were recomputed and tied to immutable
repository-object provenance, the pinned HARNESS bytes matched, the immutable binding and five exact-head workflows were
verified, 22/22 package-content hashes and lengths matched, and the 11 focused contract tests passed. The reviewer therefore
closed N7 and N8.

\subsection*{Independence qualification}

The reviewer did not design the benchmark, prepare the package or author the reviewed protocol changes, but disclosed that
they authored several earlier reviews in the same series. Because earlier findings shaped parts of the current design, the
reviewer explicitly identified a confirmation-risk channel and recommended a second reviewer with no prior history if the
programme wants stronger independence. This is recorded as a precaution rather than as a condition that invalidates the
current verdict.

\section{What the ACCEPT means --- and what it does not}

The design is considered valid for the bounded engineering question: if executed as specified, Phase A can characterize
the engineering consequences of 512D versus 128D under exact dense top-1 fixed-work search on a bound device and workload.
The review did not require protected biometric outcomes; this confirms that engineering characterization and biometric
qualification are separable in practice for this gate.

The ACCEPT does \textbf{not} authorize platform materialization yet, because three missing records would leave surprising
results valid but insufficiently interpretable. It also does not authorize Phase A execution, biometric non-inferiority
claims, production suitability claims or any opening of SCREEN, qualification TEST, protected real-route performance or
representation geometry.

\section{Three blockers before platform materialization}

\subsection*{B1 --- cache hierarchy and achievable memory bandwidth}

The environment contract did not require enough information to explain whether a measured effect is driven mainly by
arithmetic, a cache-capacity crossing or memory bandwidth. G0 is especially diagnostic because the nominal FP32 working
sets are approximately 61.44 MB at 512D and 15.36 MB at 128D; depending on the target LLC, the two arms can fall on
different sides of a cache-capacity boundary.

The prospective correction therefore requires:
\begin{itemize}
\item L1d, L2 and L3/LLC capacities, inclusivity policy and sharing topology;
\item a named achievable-memory-bandwidth characterization method and measured value;
\item arm working-set size relative to LLC capacity;
\item achieved memory throughput relative to measured bandwidth where observable.
\end{itemize}

The platform/lock-defining choices must be frozen before bandwidth characterization. The characterization may not be used
to retune those choices under the same execution identity. If it motivates a configuration change, a new prospective
identity and the applicable review sequence are required. This preserves the accepted D3 information firewall.

\subsection*{B2 --- ratios across different saturation or censoring regimes}

A common offered QPS can place the 512D and 128D arms in different queueing regimes. A ratio between a saturated arm and an
unsaturated arm does not isolate the same quantity as a within-regime pair; likewise a right-censored recovery value cannot
be reduced to a finite ratio against a finite observation.

The correction therefore requires every arm to carry an explicit saturation regime. Quantitative paired deltas or ratios
are allowed only when both arms are in the same non-censored regime. Cross-regime and censored pairs remain reported as
observations, but receive explicit \code{NO\_COMBINED\_RATIO} status rather than a misleading scalar comparison.

\subsection*{B3 --- backend dispatch observability}

A fixed BLAS/runtime configuration can dispatch different microkernels or blocking strategies at inner dimensions 512 and
128. That difference is a legitimate consequence of dimension and is part of the measured engineering effect; it must not
be forced away. It must, however, be observable enough to diagnose an unexpected delta.

The prospective correction therefore records, per arm, the backend-reported kernel/core identifier when available, the
observability method, exact GEMV or batched-GEMM call shape, batch size, memory order, leading dimension, blocking details
when exposed and raw verbose/trace evidence. If the backend does not expose an item, the protocol records
\code{NOT\_EXPOSED\_BY\_BACKEND} rather than inventing it, and mechanism-specific microkernel claims remain
\code{NOT\_DEMONSTRATED}.

\section{Accepted limitations and future sensitivity}

Two reporting limitations are now prospective guards rather than open blockers. Every joules-per-identification value must
name offered QPS and saturation regime; this prevents a load-conditioned energy number from being quoted as a context-free
device property. The report preamble must also explain designed percentile gaps created by the frozen minimum sample
support: \code{INSUFFICIENT\_SAMPLES\_NO\_PERCENTILE\_CLAIM} is a designed result, not missing data.

A separate future sensitivity remains non-blocking for Phase A: G0 and G2 differ strongly in resident working set, so their
two-point result cannot establish intermediate behavior. G1 or another intermediate point may be added prospectively in a
later phase; absent such a measurement, the Phase A report must not interpolate between G0 and G2.

\section{Current governance state}

\begin{tabularx}{\linewidth}{>{\bfseries}l X}
\toprule
Item & State \\
\midrule
Construct validity & accepted with limitations on \code{b5720dd7...} \\
N7 / N8 & closed by construct-validity transport verification \\
B1 / B2 / B3 & prospectively corrected; focused confirmation required \\
L1 / L3 & prospective reporting guards adopted \\
L2 & future sensitivity; not a Phase A blocker \\
Platform materialization & prohibited until B1--B3 focused confirmation \\
Canonical Phase A & not admissible; not executed \\
SCREEN / TEST / protected performance / geometry & sealed \\
S4N1 / S4N2 & \code{CLOSED\_NEGATIVE} / \code{CLOSED\_NEGATIVE} \\
S4N3 & not launched \\
\bottomrule
\end{tabularx}

\section{Next admissible sequence}

\begin{enumerate}
\item preserve the construct-validity verdict as bound to \code{b5720dd7...};
\item land B1--B3 prospectively in a separate correction artifact without rewriting historical reviewed bytes;
\item run bounded technical assurance on the exact correction head;
\item obtain focused independent confirmation that B1--B3 are closed, preferably from a reviewer with no prior history
in the review series;
\item only after that confirmation, materialize target hardware, external power meter, off-device load generator,
configuration-choice provenance/attestation and exact generator identity;
\item freeze an execution-specific materialized lock, rerun exact-head assurances and complete any review required by the
provenance/lock mechanism;
\item only then may canonical Phase A become eligible for execution.
\end{enumerate}

No step authorizes opening protected Study 1B biometric outcomes.

\section{Authoritative artifacts for this transition}

\begin{itemize}
\item \path{protocol/reviews/PR50_B5720DD7_PHASE_A_CONSTRUCT_VALIDITY_REVIEW_VERDICT_2026-09-16.md}
\item \path{protocol/chronicle/STUDY1B_PR50_B5720DD7_PHASE_A_CONSTRUCT_VALIDITY_ACCEPT_WITH_LIMITATIONS_2026-09-16.yaml}
\item \path{protocol/benchmarks/STUDY1B_PHASE_A_CONSTRUCT_VALIDITY_CORRECTIONS_V0_1_2026-09-16.yaml}
\item \path{protocol/reviews/STUDY1B_PHASE_A_CONSTRUCT_VALIDITY_LIMITATIONS_CLOSURE_REVIEW_REQUEST_V0_1_2026-09-16.md}
\end{itemize}

Pinned reusable method authority:\\
\url{https://github.com/gharbonnier78/scientific-research-harness/blob/3b109adcdd9a8cba4df029d3803ee0e5cb5bdf98/HARNESS.md}

\vfill
\begin{center}
\small\textit{This addendum is explanatory and status-synchronizing. Frozen protocols, independent review verdicts and
append-only Chronicle entries remain authoritative.}
\end{center}
\end{document}
